% 06-boundary
\section{Context boundaries: from the page to the slot}\label{sec:boundary}

The profile of \cref{sec:cascade} describes a URL. An ad does not see a URL.
It sits at one spot on a page, and the reader sees it together with whatever
is on screen at that moment. An interscroller opens in the middle of an
article, right after the paragraphs the reader has just read. A sidebar banner
sits next to a ``most read'' list that may have nothing to do with the
article. A full-screen interstitial covers the page and has no position in it
at all. So we ask where the context of an ad stops. Our answer is that it
stops in three places, each with its own job, and that the ad format decides
how much each one counts (\cref{fig:boundary-phone}).

\begin{figure}[t]
\centering
\begin{tikzpicture}[
  ln/.style={line width=0.9pt,draw=ink!35,line cap=round},
  big/.style={line width=1.5pt,draw=ink!40,line cap=round},
  hatch/.style={draw=ink!55,dashed,pattern=north east lines,pattern color=ink!30},
  lab/.style={font=\footnotesize,text=ink,align=right,anchor=east,inner sep=1pt},
  rlab/.style={font=\footnotesize,text=ink,align=left,anchor=west,inner sep=1pt},
  lead/.style={draw=ink!55,thin}]
  % Left: the whole page, to scale vertically (1 viewport = 1.05 cm).
  \begin{scope}[yscale=-1.05]
    \fill[soft] (-0.62,-3.85) rectangle (1.47,3.2);
    \draw[ink!60,densely dotted,thick,rounded corners=2pt] (-0.62,-3.85) rectangle (1.47,3.2);
    \fill[accent!6] (-0.42,-3) rectangle (1.27,3);
    \draw[accent!70,dashed,thick,rounded corners=2pt] (-0.42,-3) rectangle (1.27,3);
    \fill[accent!16] (-0.22,-1) rectangle (1.07,1);
    \draw[accent,thick,rounded corners=2pt] (-0.22,-1) rectangle (1.07,1);
    \fill[white] (0,-4.35) rectangle (0.85,3.8);
    \draw[ink!45] (0,-4.35) rectangle (0.85,3.8);
    \fill[hatch] (0.06,-4.25) rectangle (0.79,-4.02);
    \draw[ink!60,line width=2pt] (0.1,-3.68) -- (0.65,-3.68);
    \foreach \y in {-3.48,-3.36,...,-2.9} \draw[ln] (0.1,\y) -- (0.75,\y);
    \fill[ink!15] (0.1,-2.75) rectangle (0.75,-2.15);
    \foreach \y in {-2.0,-1.88,...,-1.2} \draw[ln] (0.1,\y) -- (0.75,\y);
    \draw[big] (0.1,-1.0) -- (0.5,-1.0);
    \foreach \y in {-0.85,-0.73,...,-0.1} \draw[ln] (0.1,\y) -- (0.75,\y);
    \fill[accent!85] (0.14,0) rectangle (0.71,0.30);
    \draw[ln] (0.1,0.35) -- (0.75,0.35);
    \fill[hatch] (0.1,0.41) rectangle (0.75,0.58);
    \foreach \y in {0.70,0.82,...,1.95} \draw[ln] (0.1,\y) -- (0.75,\y);
    \draw[big] (0.1,2.12) -- (0.45,2.12);
    \foreach \y in {2.27,2.39,...,2.8} \draw[ln] (0.1,\y) -- (0.75,\y);
    \fill[hatch] (0.06,3.3) rectangle (0.79,3.7);
    % The viewport.
    \draw[ink,line width=1.3pt,rounded corners=1.5pt] (-0.06,-0.40) rectangle (0.91,0.60);
    \coordinate (vt) at (0.91,-0.40);
    \coordinate (vb) at (0.91,0.60);
    % Labels.
    \node[lab] (a) at (-0.95,-3.45) {page: the story the slot\\sits in, with weight $w_f$};
    \draw[lead] (a.east) -- (-0.62,-3.45);
    \node[lab] (b) at (-0.95,-2.0) {local window, $|d|\le 3$:\\the 12 nearest blocks,\\weighted by distance};
    \draw[lead] (b.east) -- (-0.42,-2.0);
    \node[lab] (c) at (-0.95,-0.72) {co-visible window, $|d|\le 1$:\\its worst block decides suitability};
    \draw[lead] (c.east) -- (-0.22,-0.72);
    \node[lab] (e) at (-0.95,0.25) {viewport: on screen now};
    \draw[lead] (e.east) -- (-0.06,0.25);
    \node[lab] (f) at (-0.95,3.5) {navigation and comments\\lie outside the story};
    \draw[lead] (f.east) -- (0,3.5);
  \end{scope}
  % Right: the viewport, enlarged, as the reader's phone shows it.
  \fill[ink!5] (vt) -- (3.05,2.85) -- (3.05,-2.85) -- (vb) -- cycle;
  \draw[ink!30] (vt) -- (3.05,2.85) (vb) -- (3.05,-2.85);
  \draw[ink,line width=1.4pt,fill=white,rounded corners=9pt] (2.9,-3.3) rectangle (5.9,3.3);
  \fill[ink] (4.05,3.06) rectangle (4.75,3.12);
  \draw[ink!50,line width=1.2pt,line cap=round] (4.0,-3.1) -- (4.8,-3.1);
  \draw[ink!50] (3.05,-2.85) rectangle (5.75,2.85);
  \foreach \y in {2.62,2.42,...,1.0} \draw[ln] (3.25,\y) -- (5.55,\y);
  \draw[ln] (3.25,0.82) -- (4.6,0.82);
  \fill[accent!85] (3.36,0.57) rectangle (5.44,-1.14);
  \node[font=\footnotesize\bfseries,text=white] at (4.4,-0.29) {ad slot};
  \draw[ln] (3.25,-1.32) -- (5.55,-1.32);
  \draw[ln] (3.25,-1.50) -- (5.0,-1.50);
  \fill[hatch] (3.2,-1.68) rectangle (5.6,-2.74);
  \node[font=\tiny\bfseries,anchor=north west,inner sep=1pt,fill=white] at (3.3,-1.76) {Most read};
  \node[font=\tiny,anchor=north west,align=left,text width=2.25cm,inner sep=1pt,fill=white] at (3.3,-2.02) {1.\ Man arrested after\\\phantom{1.\ }fatal stabbing};
  % Labels on the right.
  \node[rlab] (g) at (6.15,1.75) {paragraphs just read:\\they count most};
  \draw[lead] (g.west) -- (5.55,1.75);
  \node[rlab] (h) at (6.15,-0.29) {the slot, $d = 0$};
  \draw[lead] (h.west) -- (5.44,-0.29);
  \node[rlab] (i) at (6.15,-2.2) {widget, not main:\\no weight for relevance,\\but it blocks a brand-safe ad};
  \draw[lead] (i.west) -- (5.6,-2.2);
\end{tikzpicture}
\caption{Three boundaries around one slot. Left: the whole page, drawn to scale
  in viewport heights, with offsets $d$ negative above the slot. The page is
  the story, which leaves out navigation and comments (dotted). The local
  window (dashed) holds the blocks the page module reports. The co-visible
  window (solid) holds what can be on screen with the ad, widgets included.
  Right: the viewport enlarged, as the phone shows it. The widget under the ad
  has no say in what the slot is about, and the final say on whether the slot
  is safe.}
\label{fig:boundary-phone}
\end{figure}

\subsection{Three scales, two rules}\label{sec:boundary-scales}

The \emph{page} is the story that contains the slot. It says why the reader came
and what frame of mind they are in. Most of the time the story is the URL, but
live blogs, infinite scroll and multi-story pages put several stories on one
URL. The extractor records the innermost \code{<article>} of every block, and a
story with at least 400 characters of main text gets its own labels. A slot
inside that story takes them as its page.

The \emph{local window} is the main-content blocks around the slot, weighted
down with distance. It says what the reader has just read. We weight content
above the slot more than content below it, because the reader has scrolled
past what is above and has not yet seen what is below, and content can only
prime what the reader has seen. This asymmetry is a hypothesis, not a
measurement, and \cref{sec:boundary-status} says how we plan to test it.

The \emph{co-visible window} is everything on screen with the ad, including
comments, sidebars and ``most read'' lists. It is what a screenshot of the ad
would show. Boilerplate removal is right for topic and wrong for safety. A
widget headline about a stabbing is noise when we decide what the article is
about, and a real risk for the brand placed next to it. So the extractor keeps
such blocks, with the flag \code{main = false}. They never count for relevance
and always count for suitability.

Two rules follow. Relevance blends the page with the local window. Suitability
takes the worst case over everything co-visible, so a sensitive block is never diluted
by the benign text around it.

The format sets the mix. We map ad formats onto five classes by where the ad
sits against the content, not by product name (\cref{tab:boundary-formats}).
A full-screen unit has no anchor in the text, so the page dominates. A
scroll-through unit opens inside the article, so the paragraphs just read
dominate. In-article units take a middle weight, and sticky and banner units
borrow it for now. A banner's neighbours are mostly widgets, which count only
for suitability, so we expect its page weight to settle near 0.8 once fitted.

\begin{table}[t]
\centering
\caption{Format classes and their geometry (\code{Slot.geometry} in
  \code{src/lib/slot.ml}). $w_f$ is the page's weight against the local
  window; $r^{\uparrow}_f$ and $r^{\downarrow}_f$ are the kernel's reach above
  and below the slot, in viewports. All values are priors, not fits.}
\label{tab:boundary-formats}
\small
\begin{tabularx}{\linewidth}{@{}l >{\raggedright\arraybackslash}X >{\raggedright\arraybackslash}X c c c@{}}
\toprule
Class & Examples & Position against the content & $w_f$ & $r^{\uparrow}_f$ & $r^{\downarrow}_f$\\
\midrule
\code{fullscreen}     & interstitials, takeovers & covers the page; no anchor in the text & 0.9 & 1 & 0.5\\
\code{scroll-through} & interscrollers, scroll reveals & opens inside the article as the reader scrolls & 0.4 & 1 & 0.5\\
\code{in-article}     & MPUs, in-read video, native & between paragraphs & 0.6 & 1 & 0.5\\
\code{sticky}         & anchors, overscrollers & on screen while the text scrolls under it & 0.6 & 1 & 0.5\\
\code{banner}         & leaderboards, skyscrapers & beside widgets, outside the article & 0.6 & 1 & 0.5\\
\bottomrule
\end{tabularx}
\par\smallskip
\raggedright\footnotesize Every class shares the co-visible window $|d| \le 1$.
\code{sticky} and \code{banner} reuse the \code{in-article} prior until they
are fitted.
\end{table}

We keep these weights as parameters because the evidence on congruence is
mixed. Congruent magazine ads were better remembered in one study
\citep{moorman2002effects}, while incongruent placements were sometimes
processed more in another \citep{dahlen2008could}. Contextual match and
obtrusive formats each lifted purchase intent, but not together
\citep{goldfarb2011online}. That last result bears directly on interscrollers
and interstitials.

\subsection{The slot's context}\label{sec:boundary-model}

Offsets are measured in viewport heights from the slot, because a viewport is
what a reader sees at once on any device. Let $d_b$ be the offset of block $b$:
negative above the slot (already read), positive below (not yet seen). A
format $f$ has the geometry $(w_f, r^{\uparrow}_f, r^{\downarrow}_f)$ of
\cref{tab:boundary-formats}, and its kernel is
\begin{equation}\label{eq:boundary-kernel}
k_f(d) =
\begin{cases}
  \exp\bigl(-|d| / r^{\uparrow}_f\bigr) & d < 0,\\
  \exp\bigl(-|d| / r^{\downarrow}_f\bigr) & d \ge 0.
\end{cases}
\end{equation}
With today's reaches a block one viewport above the slot weighs $0.37$ and a
block one viewport below weighs $0.14$. Let $p(v)$ be the page's (the story's)
score for node $v$, $s_b(v)$ the score of block $b$ ($0$ when $b$ lacks the
label), $M$ the main blocks near the slot and $C = \{b : |d_b| \le 1\}$ the
co-visible ones, main or not. For every $v$ on the page or on a block of $M$,
the slot's relevance is
\begin{equation}\label{eq:boundary-relevance}
\operatorname{rel}_f(v) = w_f\, p(v) + (1 - w_f)\,
  \frac{\sum_{b \in M} k_f(d_b)\, s_b(v)}{\sum_{b \in M} k_f(d_b)},
\end{equation}
and $\operatorname{rel}_f = p$ when $M$ is empty. Let $\mathcal{S}$ be the
Sensitive Topics branch of the taxonomy: its root and its 11 children. For
$v \in \mathcal{S}$ the slot's score is the worst case,
\begin{equation}\label{eq:boundary-suitability}
\operatorname{suit}(v) = \max\Bigl(p(v),\ \max_{b \in C} s_b(v)\Bigr),
\end{equation}
which replaces $\operatorname{rel}_f(v)$. The result has the same shape as a
page profile, so the matcher of \cref{sec:serving} runs on it unchanged.

\begin{proposition}\label{prop:boundary-monotone}
Adding a block to the slot's surroundings never lowers $\operatorname{suit}(v)$
for any $v \in \mathcal{S}$. Adding a main block can lower
$\operatorname{rel}_f(v)$.
\end{proposition}
\begin{proof}
A maximum over a larger set is at least the maximum over a smaller one. For
relevance, a main block $b'$ with $s_{b'}(v) = 0$ adds $k_f(d_{b'}) > 0$ to the
denominator of \cref{eq:boundary-relevance} and nothing to its numerator.
\end{proof}

The asymmetry is the point. Benign text near the slot dilutes relevance, as it
should, but it cannot hide a sensitive block. It also explains how unknown
blocks are treated at bid time (\cref{sec:boundary-where}): a block the
profile does not know counts as main content without labels, so it dilutes the
neighbourhood and never makes the slot safer.

\Cref{tab:boundary-example} shows the effect on the test that pins
\code{Slot.context}. A weekend food guide has page labels that lean towards
wine. The slot sits 0.3 viewport below the recipe section. As a scroll-through
unit, the recipe wins and cookware ranks first. As a full-screen unit, the page
wins and beer ranks first. A ``most read'' widget with a crime headline,
0.4 viewport below the slot, blocks both creatives for brand suitability. Move
that widget 2.5 viewports away, off screen, and the block lifts. The test
asserts all of these cases.

\begin{table}[t]
\centering
\caption{One page, two creatives, four slot contexts (\code{test\_contextual.ml}).
  The page scores Food \& Drink 0.85, Cooking 0.5 and Alcoholic Beverages 0.85.
  Main blocks sit at $-0.3$ (Food \& Drink 0.9, Cooking 0.9), $-2.0$ and $+1.5$
  viewports (Food \& Drink 0.9, Alcoholic Beverages 0.9). Each cell gives the
  matcher's score and its best topic pair; bold marks the winner.}
\label{tab:boundary-example}
\small
\begin{tabularx}{\linewidth}{@{}>{\raggedright\arraybackslash}p{3.6cm} >{\raggedright\arraybackslash}X >{\raggedright\arraybackslash}X@{}}
\toprule
Slot context & cookware (targets Food \& Drink) & beer\\
\midrule
\code{scroll-through}, just after the recipe & \textbf{0.341}, Cooking $\sim$ Cooking (0.57) & 0.317, Food \& Drink $\sim$ Alcoholic Beverages (0.53)\\
\code{fullscreen} & 0.308, Food \& Drink $\sim$ Cooking (0.51) & \textbf{0.423}, Alcoholic Beverages $\sim$ Alcoholic Beverages (0.70)\\
\code{scroll-through}, crime widget at $+0.4$ & blocked: brand suitability & blocked: brand suitability\\
\code{scroll-through}, crime widget at $+2.5$ & \textbf{0.341}, as in row 1 & 0.317, as in row 1\\
\bottomrule
\end{tabularx}
\end{table}

\subsection{Where the slot is}\label{sec:boundary-where}

Rendering every page in a headless browser costs 10 to 50 times a fetch, so the
crawler does not render. It finds slots in three cheaper ways.

\emph{From the template.} An ad unit's Global Placement ID (GPID), which
Prebid's GPT Pre-Auction module sets in \code{imp.ext.gpid}, usually ends with
the element id of the ad's container \citep{prebid2026global}. The extractor
records containers whose id holds a word such as \code{ad}, \code{mpu},
\code{leaderboard} or \code{gpt}, or a local-language equivalent, between
dashes, underscores or digits: \code{div-gpt-ad-1699} and
\code{top-leaderboard} match, \code{masthead} and \code{consent-banner} do not.
Each container is placed by the number of blocks before it. On the 338 pages
archived on the Pi cluster, 119 carry at least one container, once we fixed
three faults that had counted stylesheet ids, hexadecimal ids and consent
banners as slots.
Some slots are injected by script and never appear in the HTML. For those we
learn positions per template, a template being the site, its top section and
the URL depth capped at 4: a container seen on at least 2 pages of a template
gets its mean position, counted in main blocks. In-article slots usually sit at
fixed template positions, so this costs one lookup per page.

\emph{At auction time.} When the auction starts, the browser has already laid
out the page, so the best measurement is free. A Prebid.js real-time data
module (\code{sdks/prebid}) finds each ad unit's element, from the GPID's div id
or else the ad unit's code. It collects the text blocks outside the slot with
at least 20 characters of their own text, keeps those within 3 viewports, and
writes the 12 nearest into \code{ortb2Imp.ext.data.ctx}, each as a 64-bit
SimHash of its words and an offset in viewports (\cref{fig:boundary-flow}).
That is under 350 bytes of JSON per slot, and no page text leaves the page. The
module's SimHash is the crawler's, byte for byte, pinned by fixtures in five
languages that the OCaml side writes. The bidder matches each hash to the
nearest block of the page's served profile within 3 differing bits, and the
labels computed offline come with it. \Cref{app:slot} gives the wire format
and the matching rule. Prebid.js adapters built on \code{ortbConverter} copy
\code{ortb2Imp} into the OpenRTB \code{imp} unchanged. We checked this on one
such adapter on 2026-11-11. Whether Prebid Server keeps \code{imp.ext.data}
when it rewrites \code{imp.ext} per bidder is not verified yet.

\emph{From the estimated layout.} When a request carries a GPID but no
\code{ctx}, the bidder places the slot from the profile alone
(\code{Profile.SlotOfGPID}). The served profile stores each container's
estimated position and each labelled block's estimated middle, and the bidder
takes the 12 blocks nearest the container within 3 viewports, exactly as the
page module would. The estimate comes from text length on a 390-pixel-wide
phone with a 700-pixel viewport. It is rough, but on a single-column mobile
page the order of the blocks is close to their visual order. The lab's page
view measures the real layout on the archived page and scores slots on it
instead.

\begin{figure}[t]
\centering
\begin{tikzpicture}[
  fstep/.style={box,text width=3.05cm,minimum height=1.45cm,font=\footnotesize},
  lane/.style={font=\footnotesize\bfseries,text=ink!75,rotate=90,anchor=south}]
  % Offline row: the crawler.
  \node[fstep] (html) at (1.75,0) {page HTML, cut into\\blocks, each with\\a \code{main} flag};
  \node[fstep] (sh) at (5.55,0) {64-bit SimHash of\\each block's words};
  \node[fstep] (lab) at (9.35,0) {\tier{1} and \tier{2} labels\\per block, and its\\estimated position};
  \node[store,minimum width=2.5cm,minimum height=1.45cm,font=\footnotesize] (st) at (13.15,0) {serving\\profile};
  \draw[arrow] (html) -- (sh);
  \draw[arrow] (sh) -- (lab);
  \draw[arrow] (lab) -- (st);
  % Online row: the browser, then the bidder.
  \node[fstep] (mod) at (1.75,-2.75) {page module: the 12\\nearest blocks within\\3 viewports};
  \node[fstep,draw=accent,thick] (wire) at (5.55,-2.75) {bid request \code{ctx}:\\hash and offset per\\block, under 350 B in all};
  \node[fstep] (match) at (9.35,-2.75) {bidder: each hash to\\the profile block\\within 3 bits};
  \node[fstep] (ctx) at (13.15,-2.75) {slot context,\\then the deal's rule};
  \draw[arrow] (mod) -- (wire);
  \draw[arrow] (wire) -- (match);
  \draw[arrow] (match) -- (ctx);
  \draw[arrow] (st.south) -- ++(0,-0.45) -| (match.north);
  \node[note,anchor=west] at ($(st.south)+(0.08,-0.22)$) {lookup};
  \draw[dashedarrow,{Stealth[length=2mm]}-{Stealth[length=2mm]}] (sh.south) -- (wire.north);
  \node[note,anchor=east] at (5.4,-1.375) {same hash, byte for byte};
  \node[note,below=1mm of wire.south] {no page text leaves the page};
  \node[lane] at (-0.2,0) {crawler};
  \node[lane] at (-0.2,-2.75) {auction};
\end{tikzpicture}
\caption{From the page module to the slot's context. Offline, the crawler labels
  each block and stores its SimHash in the serving profile. At auction time,
  the browser sends only hashes and offsets. The bidder joins them to the
  profile's blocks and computes the slot's context by
  \cref{eq:boundary-relevance,eq:boundary-suitability} with the format's
  geometry.}
\label{fig:boundary-flow}
\end{figure}

\subsection{Slot odds are priors}\label{sec:boundary-odds}

The lab ranks slots for a creative described in words. Each slot gets two
numbers. Fit is the chance that the creative suits the slot, given that the
slot is seen. Reach is the chance that the reader scrolls as far as the slot.
Their product is the slot's odds:
\begin{equation}\label{eq:boundary-odds}
\text{fit} = \sigma\bigl(12\,(\text{score} - 0.25)\bigr),\qquad
\text{reach} = e^{-y/6},\qquad
P = \text{fit} \times \text{reach},
\end{equation}
where score is the matcher's score on the slot's context (\cref{sec:serving}),
fit is $0$ when a hard rule blocks the creative, and $y$ is the slot's depth
in viewports from the top of the page. A slot 4.2 viewports down with a fit of
0.88 gets a reach of 0.50 and odds of 0.44. Slots go where the page's
template puts its ad containers. Without any container, a slot is tried after
each main paragraph, quote, caption, image or embed. A full-screen format gets
a single page-level slot.

\begin{figure}[t]
\centering
\begin{tikzpicture}
\begin{groupplot}[group style={group size=3 by 1,horizontal sep=1.55cm},
  scale only axis,width=3.85cm,height=3.0cm,
  tick label style={font=\scriptsize},label style={font=\footnotesize},
  title style={font=\footnotesize},axis line style={ink!70},
  every axis plot/.append style={thick},grid=major,grid style={ink!10}]
\nextgroupplot[title={fit against matcher score},xlabel={matcher score},
  ylabel={fit},xmin=0,xmax=1,ymin=0,ymax=1.05,xtick={0,0.25,0.5,0.75,1},
  xticklabels={0,0.25,0.5,0.75,1}]
  \addplot[accent,domain=0:1,samples=80] {1/(1+exp(-12*(x-0.25)))};
  \addplot[only marks,mark=*,mark size=1.8pt,ink,samples at={0.341,0.423}] {1/(1+exp(-12*(x-0.25)))};
  \node[font=\tiny,anchor=north west,align=left] at (axis cs:0.345,0.72) {cookware,\\scroll-through};
  \node[font=\tiny,anchor=west] at (axis cs:0.47,0.855) {beer, fullscreen};
  \draw[ink!50,dashed,thin] (axis cs:0.25,0) -- (axis cs:0.25,0.5) -- (axis cs:0,0.5);
\nextgroupplot[title={reach against slot depth},xlabel={depth $y$ (viewports)},
  ylabel={reach},xmin=0,xmax=15,ymin=0,ymax=1.05,xtick={0,3,6,9,12,15}]
  \addplot[accent,domain=0:15,samples=60] {exp(-x/6)};
  \addplot[only marks,mark=*,mark size=1.8pt,ink,samples at={4.2}] {exp(-x/6)};
  \node[font=\tiny,anchor=south west] at (axis cs:4.4,0.52) {$y=4.2$: reach 0.50};
\nextgroupplot[title={kernel against offset},xlabel={offset $d$ (viewports)},
  ylabel={$k_f(d)$},xmin=-3,xmax=3,ymin=0,ymax=1.05,xtick={-3,-2,-1,0,1,2,3}]
  \fill[accent!12] (axis cs:-1,0) rectangle (axis cs:1,1.05);
  \addplot[accent,domain=-3:0,samples=40] {exp(x)};
  \addplot[accent,domain=0:3,samples=40,forget plot] {exp(-2*x)};
  \node[font=\tiny,anchor=south] at (axis cs:0,0.02) {co-visible};
  \node[font=\tiny,anchor=south] at (axis cs:-2,0.16) {read};
  \node[font=\tiny,anchor=south] at (axis cs:2,0.06) {unseen};
\end{groupplot}
\end{tikzpicture}
\caption{The priors the lab uses, as plotted functions. Left: fit is a logistic
  of the matcher's score and crosses 0.5 at a score of 0.25; the dots are the
  two winners of \cref{tab:boundary-example}. Middle: reach decays with depth;
  the dot is the slot of the placement API's example. Right: the kernel of
  \cref{eq:boundary-kernel}, shared by every format today, falls twice as fast
  below the slot as above it. The shaded band is the co-visible window.}
\label{fig:boundary-priors}
\end{figure}

All of \cref{eq:boundary-odds} is hand-set. The logistic's slope and centre,
the six-viewport scroll decay, the page weights, the reaches and the width of
the co-visible window are priors chosen to be plausible, and each is marked as
a shortcut in the code. None has been fitted, because fitting needs engagement
data per slot, which we do not have yet. Fit should become the engagement
regression of \cref{sec:serving}, and reach should become measured scroll-depth
curves per site. Until then $P$ ranks slots within a page. It is not a
calibrated probability, and we do not report it as one.

\subsection{What is built and what is not}\label{sec:boundary-status}

As of 2026-11-11 the OCaml \code{Slot.context} computes
\cref{eq:boundary-relevance,eq:boundary-suitability}, and the served profile
carries up to 64 labelled blocks and the page's ad containers. The Go SDK's
\code{Profile.AtSlot} reproduces \code{Slot.context} on 2,000 fixture cases, in
1.8\,\textmu s with 23 allocations for five blocks, and
\code{Profile.SlotOfGPID} agrees with the OCaml side on 1,000. The page module
hashes as the crawler does. In the lab, the page view replays the archived
page with scripts off and draws block scores, slots, the kernel and the
co-visible window on it.

The rest is planned or unmeasured. No prior of \cref{tab:boundary-formats} or
\cref{eq:boundary-odds} is fitted, and the \code{sticky} and \code{fullscreen}
geometry has not been checked against each format's specification. The served
profile still estimates block positions. The module is not yet on publishers'
pages, and the bidder integration that would read \code{ctx} is open
(\cref{sec:serving}). Bid-time geometry also drifts from render time for
sticky and refreshing slots. We dropped three ideas for now: rendering one
page per template for desktop layouts, observing after the impression what was
co-visible, and using the previous page as context for interstitials, which
would bring in user history and consent obligations.

We plan to settle how far context reaches with data. The engagement model gets
two topic features per format, page-topic match and local-topic match, and
their fitted weights say how much nearby content matters. A script fits the
page feature at site grain on 30 days of delivery logs, because delivery is
reported per site; we report no result from it here. The local feature needs
slot geometry in the logs, which only the page module provides. The decisive
test is an A/B test of page context against slot context, crossed with format
(\cref{sec:protocol}). It has not run.
